% GENERATED FILE. Edit canonical lessons, then run npm run generate:books.
% canonical-source-hash: fnv1a64:55c8335b95dbefcd
% canonical-lessons: ML-C155-kavil, ML-C155-tati, ML-C155-teruvu, ML-C155-granthasala, ML-C155-sarttu

\chapter{Cheek, Chin, Street, Library, Shirt}
\label{ch:ml-cheek-chin-street-library-shirt}
\hlchaptermodality{155}

\begin{chapteropening}
\textbf{By the end of this chapter:} \emph{I can say a cheek, a chin, a street, a library and a shirt.}

Complete the last lesson of chapter 155: I can say a cheek, a chin, a street, a library and a shirt.
\end{chapteropening}

\section[\texorpdfstring{\ml{കവിൾ} (kavil)}{കവിൾ (kavil)}]{\ml{കവിൾ} (kavil) — a cheek}

\label{lesson:ML-C155-kavil}



\begin{quote}
Before the new one: say the Malayalam for skin, then the Malayalam for a forehead.
\end{quote}

\subsection*{You'll want to know: \ml{കവിൾ}}
\textbf{\ml{കവിൾ}} — \emph{kavil} — \textquotedblleft{}a cheek\textquotedblright{}.

\begin{grammarlens}[title={What you've built}]
One more word for this chapter.
\end{grammarlens}

\subsection*{Your turn}
\begin{itemize}
\raggedright
  \item \emph{Say it:} \emph{kavil}
  \item \emph{Say it:} \emph{kavil}, once more
  \item \emph{Say it:} say \emph{kavil}
  \item \emph{Recall:} say the Malayalam for skin, then the Malayalam for a forehead, then say \emph{kavil} again
\end{itemize}

\begin{tcolorbox}[breakable,colback=teal!4,colframe=teal!35!black,title={Before you move on}]
What does \ml{കവിൾ} mean? (\textquotedblleft{}A cheek\textquotedblright{}.) Say it once more.
\end{tcolorbox}
\section[\texorpdfstring{\ml{താടി} (tāṭi)}{താടി (tāṭi)}]{\ml{താടി} (tāṭi) — a chin, a beard}

\label{lesson:ML-C155-tati}



\begin{quote}
Before the new one: say the Malayalam for a forehead, then the Malayalam for a cheek.

One frame, three meanings, all after \emph{enikkŭ}: which ending is \textquotedblleft{}I like\textquotedblright{}? (\emph{Iṣṭamāṇŭ}.)
\end{quote}

\subsection*{You'll want to know: \ml{താടി}}
\textbf{\ml{താടി}} — \emph{tāṭi} — \textquotedblleft{}a chin, a beard\textquotedblright{}.

\begin{grammarlens}[title={What you've built}]
One more word for this chapter.
\end{grammarlens}

\subsection*{Your turn}
\begin{itemize}
\raggedright
  \item \emph{Say it:} \emph{tāṭi}
  \item \emph{Say it:} \emph{tāṭi}, once more
  \item \emph{Say it:} say \emph{tāṭi}
  \item \emph{Recall:} say the Malayalam for a forehead, then the Malayalam for a cheek, then say \emph{tāṭi} again
\end{itemize}

\begin{tcolorbox}[breakable,colback=teal!4,colframe=teal!35!black,title={Before you move on}]
What does \ml{താടി} mean? (\textquotedblleft{}A chin, a beard\textquotedblright{}.) Say it once more.
\end{tcolorbox}
\section[\texorpdfstring{\ml{തെരുവ്} (teruvŭ)}{തെരുവ് (teruvŭ)}]{\ml{തെരുവ്} (teruvŭ) — a street}

\label{lesson:ML-C155-teruvu}



\begin{quote}
Before the new one: say the Malayalam for a cheek, then the Malayalam for a chin.
\end{quote}

\subsection*{You'll want to know: \ml{തെരുവ്}}
\textbf{\ml{തെരുവ്}} — \emph{teruvŭ} — \textquotedblleft{}a street\textquotedblright{}.

\begin{grammarlens}[title={What you've built}]
One more word for this chapter.
\end{grammarlens}

\subsection*{Your turn}
\begin{itemize}
\raggedright
  \item \emph{Say it:} \emph{teruvŭ}
  \item \emph{Say it:} \emph{teruvŭ}, once more
  \item \emph{Say it:} say \emph{teruvŭ}
  \item \emph{Recall:} say the Malayalam for a cheek, then the Malayalam for a chin, then say \emph{teruvŭ} again
\end{itemize}

\begin{tcolorbox}[breakable,colback=teal!4,colframe=teal!35!black,title={Before you move on}]
What does \ml{തെരുവ്} mean? (\textquotedblleft{}A street\textquotedblright{}.) Say it once more.
\end{tcolorbox}
\section[\texorpdfstring{\ml{ഗ്രന്ഥശാല} (granthaśāla)}{ഗ്രന്ഥശാല (granthaśāla)}]{\ml{ഗ്രന്ഥശാല} (granthaśāla) — a library}

\label{lesson:ML-C155-granthasala}



\begin{quote}
Before the new one: say the Malayalam for a chin, then the Malayalam for a street.

With \emph{enikkŭ} in front, say \textquotedblleft{}I know,\textquotedblright{} \textquotedblleft{}I understood\textquotedblright{} and \textquotedblleft{}I like.\textquotedblright{} (\emph{Aṟiyāṁ}, \emph{manassilāyi}, \emph{iṣṭamāṇŭ}.)
\end{quote}

\subsection*{You'll want to know: \ml{ഗ്രന്ഥശാല}}
\textbf{\ml{ഗ്രന്ഥശാല}} — \emph{granthaśāla} — \textquotedblleft{}a library\textquotedblright{}.

\begin{grammarlens}[title={What you've built}]
One more word for this chapter.
\end{grammarlens}

\subsection*{Your turn}
\begin{itemize}
\raggedright
  \item \emph{Say it:} \emph{granthaśāla}
  \item \emph{Say it:} \emph{granthaśāla}, once more
  \item \emph{Say it:} say \emph{granthaśāla}
  \item \emph{Recall:} say the Malayalam for a chin, then the Malayalam for a street, then say \emph{granthaśāla} again
\end{itemize}

\begin{tcolorbox}[breakable,colback=teal!4,colframe=teal!35!black,title={Before you move on}]
What does \ml{ഗ്രന്ഥശാല} mean? (\textquotedblleft{}A library\textquotedblright{}.) Say it once more.
\end{tcolorbox}
\section[\texorpdfstring{\ml{ഷർട്ട്} (ṣarṭṭŭ)}{ഷർട്ട് (ṣarṭṭŭ)}]{\ml{ഷർട്ട്} (ṣarṭṭŭ) — a shirt}

\label{lesson:ML-C155-sarttu}



\begin{quote}
Before the new one: say the Malayalam for a street, then the Malayalam for a library.
\end{quote}

\subsection*{You'll want to know: \ml{ഷർട്ട്}}
\textbf{\ml{ഷർട്ട്}} — \emph{ṣarṭṭŭ} — \textquotedblleft{}a shirt\textquotedblright{}.

\begin{grammarlens}[title={What you've built}]
One more word for this chapter.
\end{grammarlens}

\subsection*{Your turn}
\begin{itemize}
\raggedright
  \item \emph{Say it:} \emph{ṣarṭṭŭ}
  \item \emph{Say it:} \emph{ṣarṭṭŭ}, once more
  \item \emph{Say it:} say \emph{ṣarṭṭŭ}
  \item \emph{Recall:} say the Malayalam for a street, then the Malayalam for a library, then say \emph{ṣarṭṭŭ} again
\end{itemize}

\begin{tcolorbox}[breakable,colback=teal!4,colframe=teal!35!black,title={Before you move on}]
What does \ml{ഷർട്ട്} mean? (\textquotedblleft{}A shirt\textquotedblright{}.) Say it once more.
\end{tcolorbox}
